\section{Related work}
\label{sec:related}

\paragraph{Cheaper evaluation within one benchmark.}
Most work on evaluation cost shrinks a \emph{single} benchmark: choosing few items that preserve model rankings or scores, via clustering of model confidences \citep{vivek2024anchor}, item-response models \citep{polo2024tiny}, or a sparse, psychometrically selected item set built from thousands of models' results on six benchmarks \citep{kipnis2024metabench}; \citet{perlitz2024efficient} frame the general problem of ``efficient benchmarking'' and study how far a benchmark can be reduced. Our question is orthogonal: which \emph{whole benchmarks} to run, and what they say about the others.

\paragraph{Relations between benchmarks.}
Large suites report many tasks per model \citep{hendrycks2020mmlu,muennighoff2022mteb,liang2022helm,zhuo2024bigcodebench}, and judge-based or arena-style evaluations add further axes \citep{zheng2023judging}. Whether benchmarks agree is itself sensitive to methodological choices \citep{perlitz2024benchbench}. Observational scaling laws \citep{ruan2024observational} and Sloth \citep{polo2024sloth} exploit the low-dimensional structure of public benchmark scores to predict performance on some benchmarks from others, often as a function of model scale or family. We use no scale or family information, treat the choice of the observed benchmarks as the object of study, and ask how accurate the inference is as a function of the number of benchmarks and of \emph{which} ones.

\paragraph{Subset selection and cold-start elicitation.}
Choosing $k$ variables to predict the rest is sensor placement in a Gaussian process \citep{krause2008near}, linear subset selection \citep{das2011submodular}, and column-subset selection. Eliciting a few ratings from a new user is the cold-start problem of recommender systems, studied extensively with active learning \citep{elahi2016survey}, and computerized adaptive testing does the same for test-takers. Leaderboards differ in ways that matter: the matrix is small and dominated by one general factor, cells are missing in a model-dependent way, and the new row may be systematically better than the history. We port the Gaussian machinery, state what it implies for $k$, and test it under these conditions. The predictor side is matrix completion \citep{mazumder2010softimpute} and probabilistic PCA/factor analysis \citep{tipping1999ppca}; covariance estimation uses shrinkage \citep{ledoit2004well}.
